\documentclass[11pt,a4paper]{article}

\usepackage{geometry}
\usepackage{amsmath}
\usepackage{amssymb}
\usepackage{booktabs}
\usepackage{hyperref}

\geometry{margin=2.4cm}
\title{Spiegazione dettagliata del codice Python nella cartella \texttt{briscola/}}
\author{Progetto RL per Briscola a quattro giocatori}
\date{18 giugno 2026}

\begin{document}

\maketitle
\tableofcontents
\newpage

\section{Obiettivo del documento}

Questo documento descrive nel dettaglio i file Python presenti nella cartella
\texttt{briscola/} del progetto. L'obiettivo è chiarire:

\begin{itemize}
  \item cosa fa ogni file;
  \item quali classi, funzioni e strutture dati contiene;
  \item quali parametri possono o devono essere modificati negli esperimenti;
  \item quali parti sono già implementate;
  \item quali parti sarebbero utili da implementare in futuro.
\end{itemize}

La cartella \texttt{briscola/} contiene il nucleo logico del progetto:
regole del gioco, ambiente, osservazioni, feature, policy, baseline,
training REINFORCE, self-play, valutazione e replay. I file in \texttt{scripts/}
e \texttt{web/} usano questo nucleo, ma non sono descritti qui perché non fanno
parte della cartella richiesta.

\section{Visione d'insieme della cartella}

\begin{center}
\begin{tabular}{ll}
\toprule
File & Responsabilità principale \\
\midrule
\texttt{\_\_init\_\_.py} & Esporta le classi principali del package \\
\texttt{cards.py} & Rappresentazione di carte, mazzo, punti e ranghi \\
\texttt{rules.py} & Regole pure: squadre, presa, vincitore, punti \\
\texttt{env.py} & Ambiente completo della partita a quattro giocatori \\
\texttt{observation.py} & Osservazione legale data alla policy \\
\texttt{features.py} & Feature $\phi(M,a)$ per la policy lineare \\
\texttt{policies.py} & Policy softmax lineare e funzioni di gradiente \\
\texttt{baselines.py} & Policy random, greedy e euristica \\
\texttt{reinforce.py} & Raccolta traiettorie e update REINFORCE \\
\texttt{self\_play.py} & Pool di snapshot e training self-play \\
\texttt{evaluation.py} & Valutazione, metriche, matrice degli scontri \\
\texttt{replay.py} & Registrazione e salvataggio di partite \\
\bottomrule
\end{tabular}
\end{center}

\section{\texttt{\_\_init\_\_.py}}

\subsection{Scopo}

Il file \texttt{\_\_init\_\_.py} trasforma la cartella \texttt{briscola/} in un
package Python importabile. Serve a rendere più comodo l'uso delle classi
principali.

Invece di scrivere:

\begin{verbatim}
from briscola.env import FourPlayerBriscolaEnv
\end{verbatim}

si può scrivere:

\begin{verbatim}
from briscola import FourPlayerBriscolaEnv
\end{verbatim}

\subsection{Elementi esportati}

Il file esporta:

\begin{itemize}
  \item \texttt{Card}, \texttt{PlayedCard};
  \item \texttt{RANKS}, \texttt{SUITS}, \texttt{make\_deck};
  \item \texttt{FourPlayerBriscolaEnv};
  \item \texttt{BriscolaFeatureExtractor};
  \item \texttt{LinearSoftmaxPolicy}.
\end{itemize}

La lista \texttt{\_\_all\_\_} definisce l'interfaccia pubblica del package.
Non cambia il funzionamento interno, ma comunica quali oggetti sono pensati per
essere usati dall'esterno.

\subsection{Parametri modificabili}

Non ci sono veri parametri sperimentali in questo file. Si può modificare la
lista \texttt{\_\_all\_\_} se si vogliono esportare nuove classi o funzioni.

\subsection{Cosa implementare in futuro}

\begin{itemize}
  \item Esportare anche \texttt{RandomPolicy}, \texttt{GreedyTrickPolicy},
  \texttt{SimpleHeuristicPolicy}, \texttt{SelfPlayTrainer} ed
  \texttt{evaluate\_team\_match}, se si vuole rendere il package più comodo da
  usare in notebook e script.
  \item Aggiungere una variabile \texttt{\_\_version\_\_} per tracciare la
  versione del codice.
\end{itemize}

\section{\texttt{cards.py}}

\subsection{Scopo}

\texttt{cards.py} definisce la rappresentazione delle carte e del mazzo da 40
carte. È uno dei file più importanti perché tutte le altre parti del progetto
dipendono da una rappresentazione coerente delle carte.

Il file non contiene logica di training e non contiene informazione sullo stato
della partita. Contiene solo:

\begin{itemize}
  \item semi;
  \item ranghi;
  \item valori in punti;
  \item forza relativa delle carte;
  \item funzioni per creare e leggere carte.
\end{itemize}

\subsection{Costanti principali}

\paragraph{\texttt{SUITS}}

\begin{verbatim}
SUITS = ("cups", "coins", "clubs", "swords")
\end{verbatim}

Rappresenta i quattro semi. I nomi sono in inglese per mantenere gli
identificatori ASCII e facilmente serializzabili.

\paragraph{\texttt{RANKS}}

\begin{verbatim}
RANKS = (
    "ace", "three", "king", "knight", "jack",
    "seven", "six", "five", "four", "two",
)
\end{verbatim}

L'ordine è dal più forte al più debole nella Briscola:

\[
Asso > Tre > Re > Cavallo > Fante > Sette > Sei > Cinque > Quattro > Due.
\]

\paragraph{\texttt{POINTS}}

Associa a ogni rango il valore in punti:

\[
A=11,\quad 3=10,\quad Re=4,\quad Cavallo=3,\quad Fante=2.
\]

Le altre carte valgono zero.

\paragraph{\texttt{RANK\_TO\_STRENGTH}}

Trasforma il rango in un numero usato per confrontare due carte dello stesso
seme o due briscole. Più alto è il valore, più forte è la carta.

\subsection{Classe \texttt{Card}}

\begin{verbatim}
@dataclass(frozen=True, order=True)
class Card:
    suit: str
    rank: str
\end{verbatim}

Rappresenta una carta singola. È \texttt{frozen=True}, quindi immutabile. Questo
è utile perché una carta può essere inserita in \texttt{set}, usata come chiave
di dizionario e confrontata in modo sicuro.

La classe contiene:

\begin{itemize}
  \item \texttt{\_\_post\_init\_\_}: controlla che seme e rango siano validi;
  \item \texttt{points}: restituisce il valore della carta;
  \item \texttt{strength}: restituisce la forza della carta;
  \item \texttt{id}: crea una stringa serializzabile, ad esempio
  \texttt{ace\_of\_cups};
  \item \texttt{short}: crea una forma compatta usata nei log.
\end{itemize}

\subsection{Classe \texttt{PlayedCard}}

\begin{verbatim}
@dataclass(frozen=True)
class PlayedCard:
    player_id: int
    card: Card
\end{verbatim}

Rappresenta una carta giocata pubblicamente da un certo giocatore. È usata nella
presa corrente, nella storia pubblica e nel replay.

\subsection{Funzioni}

\paragraph{\texttt{make\_deck()}}

Crea un mazzo ordinato da 40 carte:

\[
4 \text{ semi} \times 10 \text{ ranghi} = 40.
\]

\paragraph{\texttt{card\_from\_id(card\_id)}}

Ricostruisce una carta a partire da una stringa come
\texttt{three\_of\_coins}. È usata soprattutto dall'interfaccia web quando il
browser invia al server l'ID della carta scelta.

\paragraph{\texttt{total\_deck\_points()}}

Calcola la somma dei punti del mazzo. Deve restituire sempre 120.

\subsection{Parametri modificabili}

\begin{itemize}
  \item \texttt{SUITS}: si può cambiare se si vuole usare una nomenclatura
  italiana o simbolica.
  \item \texttt{RANKS}: si può cambiare solo se si modifica la variante del
  gioco. Nel progetto corrente non andrebbe cambiato.
  \item \texttt{POINTS}: si può cambiare solo per varianti non standard.
\end{itemize}

\subsection{Cosa implementare in futuro}

\begin{itemize}
  \item Aggiungere metadati grafici, per esempio simboli o immagini delle carte.
  \item Gestire più mazzi regionali mantenendo uguali le regole.
  \item Separare la rappresentazione logica della carta dalla rappresentazione
  grafica usata nell'interfaccia.
\end{itemize}

\section{\texttt{rules.py}}

\subsection{Scopo}

\texttt{rules.py} contiene le regole pure della Briscola. Per ``regole pure''
si intende codice che non dipende dallo stato nascosto dell'ambiente e non
aggiorna parametri di apprendimento. Questo rende il file facile da testare.

\subsection{Costanti}

\begin{verbatim}
TEAM_A = 0
TEAM_B = 1
NUM_PLAYERS = 4
\end{verbatim}

Il codice usa giocatori indicizzati da 0 a 3:

\[
Team A = \{0,2\}, \qquad Team B = \{1,3\}.
\]

\subsection{Funzioni sui giocatori}

\paragraph{\texttt{team\_of(player\_id)}}

Restituisce la squadra del giocatore. I giocatori pari sono in Team A, i
giocatori dispari in Team B.

\paragraph{\texttt{partner\_of(player\_id)}}

Restituisce il compagno:

\[
partner(i) = (i+2) \bmod 4.
\]

\paragraph{\texttt{left\_opponent\_of} e \texttt{right\_opponent\_of}}

Restituiscono gli avversari seduti a sinistra e a destra secondo l'ordine del
tavolo.

\paragraph{\texttt{table\_order\_from(player\_id)}}

Restituisce i quattro giocatori in ordine di tavolo partendo da un giocatore.
È usata per l'ordine di pesca dopo una presa.

\paragraph{\texttt{next\_player(player\_id)}}

Restituisce il giocatore successivo:

\[
next(i)=(i+1)\bmod 4.
\]

\subsection{Funzioni sulle prese}

\paragraph{\texttt{trick\_points(trick)}}

Somma i punti delle carte nella presa.

\paragraph{\texttt{card\_beats(challenger, incumbent, lead\_suit, trump\_suit)}}

Decide se una carta supera un'altra. La logica è:

\begin{enumerate}
  \item una briscola batte una non briscola;
  \item tra due briscole vince quella di forza maggiore;
  \item senza briscole, può vincere solo una carta del seme di apertura;
  \item tra carte del seme di apertura vince quella di forza maggiore.
\end{enumerate}

\paragraph{\texttt{trick\_winner(trick, trump\_suit)}}

Scorre tutte le carte della presa e restituisce la giocata vincente. Funziona
anche su prese parziali, quindi può essere usata dalle feature per chiedere:
``se giocassi questa carta, chi starebbe vincendo ora?''.

\subsection{Parametri modificabili}

\begin{itemize}
  \item \texttt{NUM\_PLAYERS}: nel progetto è fissato a 4. Cambiarlo richiede
  modifiche profonde in ambiente, osservazioni e feature.
  \item La composizione delle squadre è codificata in \texttt{team\_of}. Se si
  volessero posti diversi, questa funzione andrebbe cambiata.
  \item L'ordine del tavolo è codificato nelle funzioni modulo 4.
\end{itemize}

\subsection{Cosa implementare in futuro}

\begin{itemize}
  \item Supportare varianti con obbligo di risposta al seme.
  \item Rendere configurabile la composizione delle squadre.
  \item Separare completamente ``regole di Briscola'' e ``regole del tavolo'',
  per poter riusare il codice anche a 2 giocatori.
\end{itemize}

\section{\texttt{observation.py}}

\subsection{Scopo}

\texttt{observation.py} definisce che cosa vede legalmente un giocatore. Questo
file è cruciale perché impedisce information leakage, cioè evita che la policy
veda carte nascoste.

\subsection{Classe \texttt{PlayerObservation}}

\texttt{PlayerObservation} è una dataclass immutabile che contiene:

\begin{itemize}
  \item identità del giocatore e del compagno;
  \item avversari;
  \item mano del giocatore;
  \item eventuale mano del compagno nella fase finale;
  \item seme di briscola;
  \item briscola scoperta, se ancora visibile;
  \item presa corrente;
  \item storia pubblica delle carte giocate;
  \item vincitori delle prese;
  \item punteggi;
  \item leader della presa;
  \item giocatore corrente;
  \item numero di carte rimaste nel mazzo;
  \item indice della presa;
  \item posizione nella presa.
\end{itemize}

\subsection{Regola sulla mano del compagno}

Nel progetto è stata aggiunta questa regola:

\begin{quote}
Quando il mazzo è vuoto e iniziano le ultime tre prese, i compagni di squadra
possono vedere reciprocamente le proprie carte.
\end{quote}

Per questo l'osservazione contiene:

\begin{itemize}
  \item \texttt{partner\_hand\_visible};
  \item \texttt{partner\_hand}.
\end{itemize}

Prima della fase finale:

\begin{verbatim}
partner_hand_visible = False
partner_hand = ()
\end{verbatim}

Nella fase finale:

\begin{verbatim}
partner_hand_visible = True
partner_hand = tuple(carte_del_compagno)
\end{verbatim}

\subsection{Proprietà utili}

\begin{itemize}
  \item \texttt{legal\_actions}: coincide con la mano del giocatore;
  \item \texttt{score\_difference}: punteggio squadra meno punteggio avversario;
  \item \texttt{cards\_played\_count}: numero di carte pubblicamente giocate;
  \item \texttt{opponents}: coppia degli avversari.
\end{itemize}

\subsection{Parametri modificabili}

Non ci sono iperparametri numerici. Si può però decidere quali campi rendere
visibili alla policy. Questa è una scelta metodologica importante.

\subsection{Cosa implementare in futuro}

\begin{itemize}
  \item Aggiungere una rappresentazione della sequenza completa delle prese.
  \item Aggiungere un belief state approssimato sulle carte nascoste.
  \item Aggiungere osservazioni diverse per esperimenti di ablation, ad esempio
  una versione senza memoria delle carte giocate.
\end{itemize}

\section{\texttt{env.py}}

\subsection{Scopo}

\texttt{env.py} implementa l'ambiente di gioco completo. È il file che mantiene
lo stato nascosto della partita:

\begin{itemize}
  \item mani di tutti i giocatori;
  \item mazzo residuo ordinato;
  \item briscola;
  \item presa corrente;
  \item punteggi;
  \item giocatore corrente;
  \item storia pubblica.
\end{itemize}

La policy non deve accedere direttamente a questi campi. Deve usare
\texttt{observe(player\_id)}.

\subsection{Classe \texttt{StepResult}}

È il risultato di una chiamata a \texttt{env.step(card)}. Contiene:

\begin{itemize}
  \item nuova osservazione, oppure \texttt{None} se la partita è finita;
  \item flag \texttt{done};
  \item flag \texttt{trick\_completed};
  \item vincitore della presa;
  \item punti della presa;
  \item punteggi aggiornati.
\end{itemize}

\subsection{Classe \texttt{FourPlayerBriscolaEnv}}

È l'ambiente principale. I parametri iniziali sono:

\begin{verbatim}
seed: int | None = None
start_player: int = 0
\end{verbatim}

\texttt{seed} controlla la riproducibilità della partita. \texttt{start\_player}
decide chi apre la prima presa.

\subsection{Metodo \texttt{reset}}

Inizializza una partita:

\begin{enumerate}
  \item crea il mazzo da 40 carte;
  \item mischia il mazzo usando il seed;
  \item distribuisce 3 carte a ogni giocatore;
  \item rivela una carta di briscola;
  \item mette la carta rivelata in fondo al mazzo;
  \item inizializza punteggi, storia e presa corrente.
\end{enumerate}

Dopo la distribuzione iniziale:

\[
40 - 4\cdot 3 = 28
\]

quindi restano 28 carte nel mazzo. Dopo 7 prese, il mazzo si svuota. Poi si
giocano le ultime 3 prese senza pescare. La partita termina dopo 10 prese.

\subsection{Metodo \texttt{observe}}

Costruisce una \texttt{PlayerObservation}. È il punto in cui viene applicata la
separazione tra stato interno e informazione legale.

La mano del compagno viene inclusa solo se:

\begin{verbatim}
self.partner_hand_visible()
\end{verbatim}

cioè se il mazzo è vuoto.

\subsection{Metodo \texttt{step}}

Esegue una giocata:

\begin{enumerate}
  \item controlla che la carta sia legale;
  \item rimuove la carta dalla mano del giocatore;
  \item aggiunge la carta alla presa corrente e alla storia pubblica;
  \item se la presa ha 4 carte, risolve il vincitore;
  \item assegna i punti alla squadra vincitrice;
  \item fa pescare i giocatori in ordine dal vincitore della presa;
  \item aggiorna leader e giocatore corrente;
  \item controlla se la partita è terminata.
\end{enumerate}

\subsection{Metodo \texttt{\_draw\_after\_trick}}

Dopo una presa, se il mazzo non è vuoto, pesca una carta per giocatore partendo
dal vincitore. Quando viene pescata la briscola rivelata, \texttt{revealed\_trump}
diventa \texttt{None}.

\subsection{Metodo \texttt{\_is\_terminal}}

La partita termina quando:

\begin{itemize}
  \item il mazzo è vuoto;
  \item la presa corrente è vuota;
  \item tutte le mani sono vuote;
  \item sono state giocate 10 prese.
\end{itemize}

\subsection{Metodo \texttt{assert\_invariants}}

Serve nei test. Controlla:

\begin{itemize}
  \item il totale delle carte è sempre 40;
  \item non ci sono duplicati;
  \item a fine partita i punti sommano 120;
  \item ogni giocatore ha giocato 10 carte.
\end{itemize}

\subsection{Parametri modificabili}

\begin{itemize}
  \item \texttt{seed}: da cambiare per generare partite diverse.
  \item \texttt{start\_player}: utile per valutazioni bilanciate.
  \item Regola \texttt{partner\_hand\_visible}: può essere modificata se si vuole
  rendere la mano del compagno visibile in un momento diverso.
\end{itemize}

\subsection{Cosa implementare in futuro}

\begin{itemize}
  \item Un'interfaccia stile Gymnasium: \texttt{reset}, \texttt{step},
  \texttt{action\_space}, \texttt{observation\_space}.
  \item Salvataggio completo dello stato per riprodurre una partita a metà.
  \item Supporto a più varianti: senza pesca, a 2 giocatori, con obbligo di
  risposta al seme.
  \item Validazione formale più ampia delle invarianti dopo ogni singola giocata.
\end{itemize}

\section{\texttt{features.py}}

\subsection{Scopo}

\texttt{features.py} costruisce il vettore:

\[
\phi(M_t,a)
\]

dove \(M_t\) è l'osservazione/memoria legale del giocatore e \(a\) è una carta
candidata nella mano. Questo vettore è l'input della policy softmax lineare.

La policy non riceve direttamente lo stato interno; riceve solo feature costruite
da \texttt{PlayerObservation}.

\subsection{Costanti di normalizzazione}

Il file definisce:

\begin{itemize}
  \item \texttt{MAX\_CARD\_POINTS = 11};
  \item \texttt{MAX\_TRICK\_POINTS = 44};
  \item \texttt{MAX\_DECK\_COUNT = 28};
  \item \texttt{TOTAL\_POINTS = 120};
  \item \texttt{TOTAL\_TRUMPS = 10}.
\end{itemize}

Servono per portare molte feature in intervalli simili, di solito tra 0 e 1.
Questo rende il learning rate più stabile.

\texttt{MAX\_TRICK\_POINTS} vale 44 perché in una presa a quattro giocatori
possono teoricamente comparire i quattro Assi, ognuno da 11 punti. Il valore
\texttt{MAX\_DECK\_COUNT} vale invece 28 perché, dopo aver distribuito le tre
carte iniziali a ciascuno dei quattro giocatori, nel mazzo restano:
\[
40 - 4 \cdot 3 = 28
\]
carte da pescare.

\subsection{Classe \texttt{BriscolaFeatureExtractor}}

La classe contiene:

\begin{itemize}
  \item \texttt{feature\_names}: lista ordinata dei nomi delle feature;
  \item \texttt{size()}: numero di feature;
  \item \texttt{extract(obs, card)}: restituisce il vettore numerico.
\end{itemize}

Il codice attuale usa 68 feature.

\subsection{Categorie di feature}

\paragraph{Feature della carta}

Descrivono la carta candidata:

\begin{itemize}
  \item punti;
  \item forza del rango;
  \item se è briscola;
  \item se è Asso, Tre, carico o carta bassa.
\end{itemize}

\paragraph{Feature della presa}

Descrivono la presa corrente:

\begin{itemize}
  \item posizione nella presa;
  \item punti presenti nella presa;
  \item se il compagno o l'avversario sta vincendo;
  \item se la carta candidata vincerebbe ora;
  \item se batterebbe il compagno;
  \item se batterebbe l'avversario.
\end{itemize}

\paragraph{Feature sui giocatori ancora da muovere}

Misurano il rischio:

\begin{itemize}
  \item quanti giocatori devono ancora agire;
  \item quanti di questi sono avversari;
  \item se il compagno deve ancora agire.
\end{itemize}

\paragraph{Feature sulla mano del compagno}

Sono attive solo nella fase finale:

\begin{itemize}
  \item \texttt{partner\_hand\_visible};
  \item punti nella mano del compagno;
  \item numero di briscole del compagno;
  \item numero di carte alte del compagno;
  \item se il compagno può vincere la presa corrente.
\end{itemize}

\paragraph{Feature di memoria pubblica}

Usano carte già osservate:

\begin{itemize}
  \item carte giocate;
  \item briscole giocate;
  \item briscole non osservate;
  \item Assi e Tre già usciti;
  \item carte alte già uscite;
  \item carte più forti dello stesso seme non osservate;
  \item briscole più forti non osservate.
\end{itemize}

\paragraph{Feature di punteggio e fase}

Includono:

\begin{itemize}
  \item punteggio della squadra;
  \item punteggio avversario;
  \item differenza punti;
  \item squadra avanti o indietro;
  \item carte nel mazzo;
  \item fase iniziale, centrale, finale;
  \item mazzo vuoto.
\end{itemize}

\paragraph{Feature di interazione}

Sono prodotti tra segnali semplici. Per esempio:

\[
\texttt{trump\_x\_trick\_points}
=
\texttt{is\_trump}\cdot \texttt{trick\_points}.
\]

Servono perché una policy lineare non può imparare da sola condizioni come:
``gioca briscola solo se la presa vale molto''.

\subsection{Parametri modificabili}

\begin{itemize}
  \item \texttt{feature\_names}: si può passare una lista custom per fare
  ablation.
  \item Costanti di normalizzazione: devono essere coerenti con il gioco.
  \item Soglie euristiche interne, ad esempio cosa viene considerato
  \texttt{high\_card} o \texttt{risky\_card}.
\end{itemize}

\subsection{Cosa implementare in futuro}

\begin{itemize}
  \item Un estrattore \texttt{NoMemoryFeatureExtractor} per l'ablation senza
  memoria.
  \item Feature basate su belief approssimato.
  \item Test automatici più dettagliati sui valori delle feature in stati noti.
  \item Feature per distinguere meglio sicurezza e rischio quando devono ancora
  giocare più avversari.
\end{itemize}

\section{\texttt{policies.py}}

\subsection{Scopo}

\texttt{policies.py} definisce l'interfaccia comune delle policy e implementa la
policy principale del progetto: una softmax lineare con action mask.

\subsection{Protocollo \texttt{Policy}}

Definisce i metodi minimi che ogni policy deve avere:

\begin{itemize}
  \item \texttt{action\_probabilities(obs)};
  \item \texttt{select\_action(obs, rng, greedy=False)}.
\end{itemize}

Questo permette di usare nello stesso ambiente policy diverse: learner, random,
greedy, euristica, snapshot.

\subsection{Classe \texttt{LinearSoftmaxPolicy}}

È la policy appresa. Per ogni azione legale \(a\):

\[
h_\theta(M,a) = \theta^\top \phi(M,a).
\]

Poi:

\[
\pi_\theta(a|M)
=
\frac{\exp(h_\theta(M,a))}
{\sum_{a'\in H}\exp(h_\theta(M,a'))}.
\]

Solo le carte nella mano sono considerate. Questo realizza l'action masking.

\subsection{Metodi principali}

\paragraph{\texttt{initialize}}

Crea una policy con pesi casuali piccoli. Parametri modificabili:

\begin{itemize}
  \item \texttt{scale}: ampiezza iniziale dei pesi;
  \item \texttt{rng}: generatore casuale;
  \item \texttt{name}: nome della policy.
\end{itemize}

\paragraph{\texttt{action\_preferences}}

Calcola \(h_\theta(M,a)\) per ogni carta legale.

\paragraph{\texttt{action\_probabilities}}

Applica la softmax. Il codice sottrae il massimo prima dell'esponenziale per
stabilità numerica:

\[
\exp(h_a - \max_b h_b).
\]

\paragraph{\texttt{select\_action}}

Se \texttt{greedy=False}, campiona secondo le probabilità. Se
\texttt{greedy=True}, sceglie la carta con probabilità massima.

\paragraph{\texttt{grad\_log\_probability}}

Calcola:

\[
\nabla_\theta \log \pi_\theta(a_t|M_t)
=
\phi(M_t,a_t)
-
\sum_{a'\in H_t}
\pi_\theta(a'|M_t)\phi(M_t,a').
\]

Questa è la formula usata da REINFORCE.

\paragraph{\texttt{apply\_gradient}}

Aggiorna i parametri:

\[
\theta \leftarrow \theta + \alpha g.
\]

Può anche tagliare il gradiente con \texttt{max\_update\_norm}.

\subsection{Parametri modificabili}

\begin{itemize}
  \item \texttt{theta}: parametri appresi.
  \item \texttt{scale}: inizializzazione.
  \item \texttt{greedy}: modalità di selezione azione.
  \item \texttt{max\_update\_norm}: clipping del gradiente.
\end{itemize}

\subsection{Cosa implementare in futuro}

\begin{itemize}
  \item Salvataggio e caricamento della policy direttamente nella classe.
  \item Temperatura della softmax:
  \[
  \pi(a|M) \propto \exp(h(M,a)/T).
  \]
  \item Policy neurale come estensione, mantenendo la stessa interfaccia.
  \item Logging delle probabilità per spiegare le decisioni all'utente.
\end{itemize}

\section{\texttt{baselines.py}}

\subsection{Scopo}

\texttt{baselines.py} contiene policy non apprese, usate come confronto e
sanity check. Sono fondamentali per capire se il learner ha imparato qualcosa.

\subsection{\texttt{RandomPolicy}}

Sceglie uniformemente tra le carte legali:

\[
\pi(a|H)=\frac{1}{|H|}.
\]

È la baseline minima.

\subsection{\texttt{GreedyTrickPolicy}}

È una policy miope: valuta solo la presa corrente. Se una carta sta vincendo la
presa, dà valore ai punti immediati. Se non può vincere, preferisce sacrificare
carte di basso valore.

Non considera:

\begin{itemize}
  \item carte future;
  \item memoria strategica;
  \item punteggio globale;
  \item cooperazione di lungo periodo.
\end{itemize}

\subsection{\texttt{SimpleHeuristicPolicy}}

È un riferimento esterno più ragionevole. Alcune regole:

\begin{itemize}
  \item se il compagno sta vincendo, carica punti;
  \item se il giocatore può vincere, valuta i punti della presa;
  \item evita di sprecare briscole su prese povere;
  \item penalizza il dono di carichi agli avversari.
\end{itemize}

\subsection{Parametri modificabili}

Nelle baseline ci sono pesi euristici come:

\begin{itemize}
  \item \texttt{0.05 * card.points};
  \item \texttt{0.5 * card.points};
  \item \texttt{2.0 * candidate\_points};
  \item penalità \texttt{5.0} per carichi dati agli avversari.
\end{itemize}

Questi numeri non sono appresi. Sono scelte manuali. Se vengono cambiati, va
documentato perché la baseline cambia.

\subsection{Cosa implementare in futuro}

\begin{itemize}
  \item Baseline più forti, ma sempre dichiarate e non usate per addestrare il
  learner principale.
  \item Baseline ``conservativa'' e ``aggressiva'' per stress test.
  \item Baseline che usa la mano del compagno nella fase finale.
\end{itemize}

\section{\texttt{reinforce.py}}

\subsection{Scopo}

\texttt{reinforce.py} implementa l'algoritmo REINFORCE con baseline.

L'update è:

\[
\theta
\leftarrow
\theta
+
\alpha
(G_t-b_t)
\nabla_\theta \log \pi_\theta(a_t|M_t).
\]

\subsection{\texttt{RewardConfig}}

Contiene:

\begin{itemize}
  \item \texttt{mode}: tipo di reward;
  \item \texttt{lambda\_margin}: peso del margine nella reward combinata.
\end{itemize}

Default:

\begin{verbatim}
mode = "combined_terminal"
lambda_margin = 0.2
\end{verbatim}

\subsection{Reward disponibili}

\paragraph{\texttt{terminal\_win}}

\[
R_T =
\begin{cases}
+1 & \text{se vince la squadra del learner},\\
0 & \text{se pareggia},\\
-1 & \text{se perde}.
\end{cases}
\]

\paragraph{\texttt{combined\_terminal}}

\[
R_T =
\operatorname{sign}(\Delta_T)
+
\lambda \frac{\Delta_T}{120}.
\]

\paragraph{\texttt{trick\_point\_dense}}

Assegna punti positivi o negativi a ogni presa in base a chi la vince.

\subsection{\texttt{collect\_episode}}

Gioca una partita completa e salva solo le decisioni del learner attivo. Il
compagno e gli avversari sono policy esterne congelate.

Per ogni decisione del learner salva:

\begin{itemize}
  \item osservazione;
  \item azione;
  \item indice temporale globale;
  \item return-to-go, calcolato dopo la fine.
\end{itemize}

\subsection{Return-to-go}

Dopo la partita:

\[
G_t = \sum_{k=t}^{T-1} R_{k+1}.
\]

Nel codice:

\begin{verbatim}
step.reward_to_go = sum(rewards[step.global_index:])
\end{verbatim}

Questo corrisponde a \(\gamma=1\).

\subsection{\texttt{reinforce\_update}}

Applica l'update su un batch di episodi. La baseline default è:

\[
b = \text{media dei return-to-go nel batch}.
\]

Poi per ogni decisione:

\[
advantage = G_t-b.
\]

Se l'advantage è positivo, l'azione viene rinforzata. Se è negativo, viene
penalizzata.

\subsection{Parametri modificabili}

\begin{itemize}
  \item \texttt{learning\_rate}: dimensione del passo di aggiornamento.
  \item \texttt{baseline}: \texttt{batch\_mean} oppure \texttt{none}.
  \item \texttt{max\_update\_norm}: clipping del gradiente.
  \item \texttt{reward\_config.mode}: reward usata.
  \item \texttt{lambda\_margin}: peso del margine nella reward combinata.
\end{itemize}

\subsection{Cosa implementare in futuro}

\begin{itemize}
  \item Baseline leave-one-out.
  \item Normalizzazione degli advantage nel batch.
  \item Critic appreso come estensione actor-critic.
  \item Entropy bonus per mantenere esplorazione.
  \item Aggiornamento con più learner nella stessa squadra, come ablation.
\end{itemize}

\section{\texttt{self\_play.py}}

\subsection{Scopo}

\texttt{self\_play.py} implementa il self-play con pool storico di snapshot.
Serve a evitare che il learner impari solo contro una singola policy.

\subsection{\texttt{Snapshot}}

Contiene:

\begin{itemize}
  \item nome;
  \item vettore \(\theta\);
  \item indice dell'update in cui è stato salvato;
  \item eventuale \texttt{evaluation\_score}, cioè uno score esterno ottenuto
  valutando la snapshot su seed fissati contro baseline o policy di test.
\end{itemize}

\subsection{\texttt{SnapshotPool}}

È il contenitore delle policy congelate. Salva solo i vettori di parametri, non
oggetti policy vivi. Quando campiona, crea una nuova \texttt{LinearSoftmaxPolicy}
con quei pesi.

\subsection{Gestione della dimensione del pool}

La versione aggiornata non assume più che le snapshot migliori siano sempre
quelle più recenti. Se il pool supera \texttt{max\_size}, il codice conserva una
combinazione di:

\begin{itemize}
  \item la snapshot iniziale;
  \item le migliori snapshot valutate, se è disponibile
  \texttt{evaluation\_score};
  \item alcune snapshot recenti;
  \item alcune snapshot più vecchie distribuite nel tempo.
\end{itemize}

In questo modo il pool rimane piccolo, ma non collassa soltanto sulle ultime
versioni della policy. Questo è importante perché nel self-play il progresso
non è necessariamente monotono: una policy recente può battere l'ultima
avversaria ma perdere contro una strategia precedente.

Il caso \texttt{max\_size = 1} è trattato come esperimento \emph{latest-only}:
il pool conserva solo l'ultima snapshot.

\paragraph{Snapshot migliori}
Il campo \texttt{evaluation\_score} non viene prodotto automaticamente dalla
partita di training. È uno score esterno, calcolato per esempio valutando la
snapshot contro random, greedy ed euristica su seed fissati. Questa separazione
evita di confondere il ritorno rumoroso del batch di training con una misura
di qualità più stabile.

\paragraph{Snapshot distribuite nel tempo}
Se rimane spazio nel pool dopo aver conservato iniziale, migliori e recenti, il
codice seleziona alcune snapshot vecchie in modo distribuito lungo gli update.
Questa scelta mantiene diversità storica e rende meno probabile dimenticare
strategie che erano utili in fasi precedenti del self-play.

\subsection{\texttt{SelfPlayConfig}}

Parametri principali:

\begin{itemize}
  \item \texttt{batch\_size};
  \item \texttt{learning\_rate};
  \item \texttt{snapshot\_interval};
  \item \texttt{max\_update\_norm};
  \item \texttt{baseline};
  \item \texttt{reward\_config}.
\end{itemize}

\subsection{\texttt{SelfPlayTrainer}}

Per ogni update:

\begin{enumerate}
  \item genera un batch di partite;
  \item campiona il posto del learner;
  \item campiona partner e avversari dal pool;
  \item raccoglie episodi;
  \item applica REINFORCE solo al learner;
  \item ogni \texttt{snapshot\_interval} salva una nuova snapshot.
\end{enumerate}

\subsection{Parametri modificabili}

\begin{itemize}
  \item \texttt{batch\_size}: più alto riduce varianza ma costa tempo.
  \item \texttt{learning\_rate}: più alto apprende più velocemente ma può
  destabilizzare.
  \item \texttt{snapshot\_interval}: controlla quanto spesso il pool cresce.
  \item \texttt{max\_size}: dimensione massima del pool.
  \item \texttt{best\_count}: quante snapshot con score migliore mantenere.
  \item \texttt{recent\_count}: quante snapshot recenti mantenere.
  \item \texttt{keep\_initial}: se mantenere sempre la policy iniziale.
  \item \texttt{evaluation\_score}: score opzionale usato per identificare
  snapshot forti.
  \item Strategia di campionamento: attualmente uniforme.
\end{itemize}

\subsection{Cosa implementare in futuro}

\begin{itemize}
  \item Campionamento non uniforme, per esempio più peso a snapshot forti o
  difficili.
  \item Valutazione più ricca delle snapshot tramite matrice degli scontri.
  \item Pool di partner distinto dal pool di avversari.
  \item Salvataggio automatico del pool su disco.
\end{itemize}

\section{\texttt{evaluation.py}}

\subsection{Scopo}

\texttt{evaluation.py} contiene funzioni per confrontare policy congelate senza
aggiornare i parametri.

\subsection{\texttt{MatchResult}}

Risultato di una partita:

\begin{itemize}
  \item punteggi finali;
  \item squadra vincitrice;
  \item differenza punti dal punto di vista del Team A.
\end{itemize}

\subsection{\texttt{EvaluationResult}}

Risultato aggregato su molte partite:

\begin{itemize}
  \item numero di partite;
  \item win rate;
  \item draw rate;
  \item loss rate;
  \item differenza media punti;
  \item standard error;
  \item intervallo di confidenza al 95\%.
\end{itemize}

\subsection{\texttt{play\_match}}

Gioca una partita usando un dizionario:

\begin{verbatim}
policies_by_player = {
    0: policy_a,
    1: policy_b,
    2: policy_a,
    3: policy_b,
}
\end{verbatim}

Non fa training. Serve solo per valutazione.

\subsection{\texttt{evaluate\_team\_match}}

Confronta due policy come squadre:

\[
policy\_a = Team A, \qquad policy\_b = Team B.
\]

Se \texttt{paired=True}, ripete la stessa partita scambiando i posti:

\[
policy\_b = Team A, \qquad policy\_a = Team B.
\]

Questo riduce l'effetto della fortuna del seed.

\subsection{\texttt{matchup\_matrix}}

Costruisce una matrice degli scontri tra molte policy. È utile per individuare:

\begin{itemize}
  \item miglioramento reale;
  \item cicli;
  \item regressioni;
  \item overfitting contro singole snapshot.
\end{itemize}

\subsection{Parametri modificabili}

\begin{itemize}
  \item lista dei seed;
  \item \texttt{paired};
  \item \texttt{greedy};
  \item insieme di policy da confrontare.
\end{itemize}

\subsection{Cosa implementare in futuro}

\begin{itemize}
  \item Test statistici più robusti.
  \item Bootstrap degli intervalli di confidenza.
  \item Salvataggio automatico in CSV.
  \item Metriche diagnostiche: uso briscole, carichi, comportamento in finale.
\end{itemize}

\section{\texttt{replay.py}}

\subsection{Scopo}

\texttt{replay.py} permette di registrare una partita e salvarla in JSON.

\subsection{Classe \texttt{Replay}}

Contiene:

\begin{itemize}
  \item seed;
  \item start player;
  \item seme di briscola;
  \item lista delle azioni;
  \item punteggi finali.
\end{itemize}

\subsection{\texttt{record\_match}}

Gioca una partita completa con policy fissate e salva, per ogni azione:

\begin{itemize}
  \item giocatore;
  \item carta;
  \item se la presa è stata completata;
  \item vincitore della presa;
  \item punti della presa;
  \item punteggi dopo la giocata.
\end{itemize}

\subsection{\texttt{save\_replay} e \texttt{load\_replay}}

Salvano e caricano il replay in formato JSON.

\subsection{Parametri modificabili}

\begin{itemize}
  \item \texttt{seed};
  \item \texttt{start\_player};
  \item \texttt{greedy}, per scegliere valutazione deterministica o stocastica.
\end{itemize}

\subsection{Cosa implementare in futuro}

\begin{itemize}
  \item Salvare anche le probabilità assegnate alle azioni.
  \item Salvare le feature al momento della decisione.
  \item Aggiungere un visualizzatore di replay.
  \item Rendere possibile riprendere una partita da un replay parziale.
\end{itemize}

\section{Parametri principali da cambiare negli esperimenti}

\begin{center}
\begin{tabular}{lll}
\toprule
Parametro & File & Significato \\
\midrule
\texttt{learning\_rate} & \texttt{reinforce.py}, \texttt{self\_play.py} & Passo di aggiornamento \\
\texttt{batch\_size} & \texttt{self\_play.py} & Partite per update \\
\texttt{snapshot\_interval} & \texttt{self\_play.py} & Ogni quanti update salvare snapshot \\
\texttt{max\_size} & \texttt{self\_play.py} & Dimensione massima del pool \\
\texttt{best\_count} & \texttt{self\_play.py} & Quante snapshot valutate come migliori conservare \\
\texttt{recent\_count} & \texttt{self\_play.py} & Quante snapshot recenti conservare \\
\texttt{pool\_eval\_games} & \texttt{train.py} & Numero di partite per assegnare score alle snapshot \\
\texttt{reward\_config.mode} & \texttt{reinforce.py} & Tipo di reward \\
\texttt{lambda\_margin} & \texttt{reinforce.py} & Peso del margine punti \\
\texttt{baseline} & \texttt{reinforce.py} & Baseline REINFORCE \\
\texttt{max\_update\_norm} & \texttt{policies.py}, \texttt{reinforce.py} & Clipping del gradiente \\
\texttt{feature\_names} & \texttt{features.py} & Feature usate dalla policy \\
\texttt{seed} & \texttt{env.py}, \texttt{evaluation.py} & Riproducibilità \\
\bottomrule
\end{tabular}
\end{center}

\section{Cose importanti ancora da implementare}

\subsection{Ablation sulla memoria}

Serve una versione alternativa dell'estrattore di feature che non usa:

\begin{itemize}
  \item carte già giocate;
  \item briscole uscite;
  \item Assi e Tre usciti;
  \item carte superiori non osservate.
\end{itemize}

Questo permetterebbe di confrontare:

\[
\text{policy con memoria} \quad \text{vs} \quad \text{policy senza memoria}.
\]

\subsection{Latest-only contro pool storico}

Il codice supporta già una modalità \emph{latest-only}: basta usare
\texttt{max\_size = 1} nel pool. In quel caso, quando viene salvata una nuova
snapshot, il pool elimina le precedenti e mantiene solo l'ultima. Questa
configurazione serve come ablation da confrontare con il pool storico completo.

\subsection{Logging diagnostico}

Mancano metriche comportamentali come:

\begin{itemize}
  \item frequenza di uso delle briscole;
  \item punti vinti quando si usa una briscola;
  \item carichi giocati quando il compagno sta vincendo;
  \item errori apparenti nella fase finale;
  \item distribuzione delle probabilità della policy.
\end{itemize}

\subsection{Salvataggio migliore dei checkpoint}

Attualmente il training salva checkpoint tramite gli script. Sarebbe utile
spostare salvataggio e caricamento dentro classi o funzioni del package. Il
checkpoint salva già, per ogni snapshot, nome, parametri, update di creazione e
eventuale \texttt{evaluation\_score}.

\subsection{Baseline più forti}

Random, greedy ed euristica sono sufficienti per iniziare, ma una valutazione
più convincente potrebbe includere:

\begin{itemize}
  \item euristica esperta;
  \item euristica aggressiva;
  \item euristica conservativa;
  \item policy addestrata con reward diversa.
\end{itemize}

\subsection{Belief state approssimato}

La policy non calcola probabilità esplicite sulle carte nascoste. Un'estensione
possibile è stimare:

\[
P(\text{compagno ha una briscola alta} \mid \text{storia pubblica}).
\]

Questo collegherebbe più direttamente il progetto alla teoria dei Bayesian
POMDP.

\subsection{Actor-critic o baseline appresa}

Il progetto usa una baseline semplice, cioè la media del batch. Una possibile
estensione è apprendere una funzione valore:

\[
V_\psi(M_t)
\]

e usare:

\[
G_t - V_\psi(M_t)
\]

come advantage. Questo trasformerebbe però il progetto verso actor-critic, che
non è il nucleo attuale.

\section{Conclusione}

La cartella \texttt{briscola/} implementa già un progetto RL completo nella sua
versione essenziale:

\begin{itemize}
  \item motore di gioco corretto;
  \item osservazioni senza leakage;
  \item policy softmax lineare;
  \item feature interpretabili;
  \item baseline;
  \item REINFORCE con baseline;
  \item self-play con pool storico;
  \item valutazione e replay.
\end{itemize}

Le estensioni più importanti da aggiungere per una relazione sperimentale forte
sono: ablation sulla memoria, confronto latest-only contro pool, logging
diagnostico, metriche più ricche e gestione più strutturata dei checkpoint.

\end{document}
